\section{Discussion}
\label{sec:discussion}

\paragraph{What the knockout shows.} Without \sig{} in pretraining or replay, 16 rounds of this EI acquired no classical
family on any of six seeds. With it, full-exposure seeds gained 8--69 strict held-out theorems during the RL phase, and a
1\% dose was enough on three seeds. The pre-registered verdict, \textsc{needs exposure}, holds for converse
contraposition, the family the full bases mostly had before RL; for the others it is undetermined at six seeds. Beyond
the waiting-time null (post hoc), the knockout ladders raised their own rate of the free positive reductio: three reached
a strict training theorem sooner than their base rates predict. They did not raise it to a family.

\paragraph{Four readings of the null.}
\begin{enumerate}
\item \emph{\sig{} is out of the policy's reach.} This is false for five of six bases, which proved strict theorems with
\sig{} at large $k$, and for four of six under guided decoding at $k = 256$. Seed 0's base is certified out of reach only
for plain sampling, at one budget and with an 8\% margin; its guided evaluation proved a strict theorem.
\item \emph{EI samples too little.} At 32 attempts per theorem, EI's per-round budget on the strict theorems is about
4,000 attempts, against base rates near $10^{-6}$. That explains why most ladders saw nothing for many rounds. It does
not explain everything: three ladders reached strict theorems sooner than this null predicts ($p = 0.006$--$0.019$).
\item \emph{EI amplifies rare finds too slowly, or too narrowly.} Seed 5 wrote a free positive reductio in round 2 and did
not spread it in 14 more rounds. Seed 1 multiplied its free positive reductio more than tenfold in eight extra rounds and still
proved the same five held-out theorems. In EI a single new proof enters a fine-tuning mix of about 70,000 RL records (up
to four proofs for each of 4,495 theorems, each repeated four times) plus 20,000 replay records.
\item \emph{The knockout made the bases worse learners in general.} The general costs are a few theorems out of
hundreds, and the knockout's gains elsewhere are not clearly smaller (partly catch-up). The dose arm restores
acquisition while moving the constituents by about one point. This reading is not excluded, but it is not needed.
\end{enumerate}

\paragraph{Positive-only or signed updates?} Whether another RL algorithm would do better from the knockout bases is the
main open question. EI updates only on successes. GRPO also pushes down failures and gives a lone success in a group of
32 an advantage of $\sqrt{31} \approx 5.6$ standard deviations. That helps only after a first success: when all 32
attempts fail, the group baseline gives no gradient, so on the strict theorems GRPO at $k=32$ faces the same waiting time
as EI. The evidence cuts both ways.
\begin{itemize}
\item \textbf{Against EI-specificity.} In this organism, from the same full-exposure bases at equal sampled attempts,
GRPO did not prove more theorems than EI. Default and unlikeliness-weighted GRPO proved significantly fewer of the 322
(theorems solved only by GRPO vs only by EI: 15 vs 38, sign test $p = 0.0022$; 15 vs 44, $p = 0.0002$). A pass@4
advantage was level (33 vs 20, $p = 0.098$), ahead on the hardest theorems mostly through one seed. Interplay's 0\%
exposure null used GRPO \citep{zhang2025interplay}. A strategy removed at fine-tuning did not re-emerge under GRPO
\citep{krishnamurthy2026select}. In LLM reasoning EI matched or beat PPO \citep{havrilla2024teaching}.
\item \textbf{For EI-specificity.} Rejection-sampling fine-tuning plateaued where GRPO built new shortcuts
\citep{abdulsalam2026rewrite}, iterated RFT failed where RL composed functions \citep{yuan2025fgx}, and GRPO produced a
pattern family absent from its RL data \citep{park2025countdown}. Training on negative samples preserves breadth that
positive-only training loses \citep{zhu2025nsr}. In our earlier 3.2M organism, GRPO acquired a depth-3 nesting habit from
zero-rate bases where EI mostly did not (below).
\end{itemize}
GRPO has not been run on the reductio in any of our organisms. A GRPO arm from the six knockout bases, with per-seed
predictions written down first, is the decisive experiment. By sampling alone, only the seeds with a strict success
inside 16 rounds (1, 2, 5) could be rescued, so any further rescue would have to come from drift.

\paragraph{Earlier organisms.} The project studied related patterns in smaller organisms: 3.2M--85M-parameter models
that write whole proofs, pretrained on proofs of at most six lines (their RL proofs ran to seven or eight). They were
judged by an older in-house checker, with Lean confirmation where stated. Their criterion for \emph{igniting} a pattern
was at least 2\% of the targets that need it, much weaker than our family bar.
\begin{itemize}
\item \textbf{Reductio, required pools.} Four of eleven reductio-free bases emitted the strict reductio before RL
($8\times10^{-6}$--$2.3\times10^{-4}$ per sample), and exactly those four ignited under EI. First proofs arrived about
when $1/(rkN)$ predicts (rounds 1, 2, 6 and 6, against 0.2, 1.7, 2.6 and 6.2). The other seven stayed at 0 of 606.
In those pools every target needed the pattern, so a zero-rate base gave EI no training step at all. At 3.2M, 25M and
85M parameters the rates did not rise with size. A rare 85M base that read 0 of 300 after eight rounds produced its
first proof in round 13 when continued. Igniting arms saturated the seven-line targets and went no further, an analogue
of knockout seed 1's narrow spread.
\item \textbf{Injected data.} One fine-tuning step on a sibling's found proofs ignited all 11 dead arms. A single injected
proof ignited the reductio on two dead seeds, reaching 6 and 17 targets in the next round. Four unrelated generator
instances, and even four invalid pattern strings, ignited some arms. In the present organism, 16 demonstrations taught
excluded middle (35--37 of 39 members) to a base pretrained with no \texttt{DN}.
\item \textbf{Mixed pools.} Where EI could train on other proofs, it ignited a depth-3 nesting habit from zero on three
of five bases, decided during training rather than by sampling luck. GRPO ignited it on every surviving arm from two
zero-rate bases (12 of 12), where paired EI did so in 1 of 3. That habit was not needed for the targets, the zero rates
were measured with the older checker, and the checkpoints are lost.
\end{itemize}
The knockout fits these results only partly. As there, the base rate predicts which seeds move, and no ladder that
started at a pooled rate near $10^{-6}$ reached the family bar. Unlike the required pools, our mixed pool let the
knockout ladders' rate rise: by the earlier criterion, knockout seed 1 ignited at round 13, and three ladders reached
strict theorems sooner than their base rates predict. ``Needs exposure'' should therefore be read as ``needs exposure,
demonstrations, or perhaps a different update, for this EI to acquire the family within 16 rounds''.

\paragraph{An amplification threshold?} One might expect EI to amplify a composition once some training theorems' rates
are high enough for 32 attempts per round to hit them within a few rounds. Our data are mixed. Knockout seed 1 amplified
members its base proved at about $10^{-4}$, but its largest gain was on a member with no base success. Knockout seed 4's
base proved three members at $1.2$--$4.9\times10^{-4}$, and its ladder amplified none. Two doses cannot tell a threshold
from a graded response. The median likelihood of known proofs also misses it: the nearest dose and knockout bases are
within 0.7 nats, yet their outcomes differ by 28--56 theorems. The 0.1\% dose, and per-theorem rates on training
theorems at several checkpoints, would test it.

\paragraph{What transfers to larger models.} The protocol transfers: budgets matched to RL's compute and reported as
curves, RL-free controls, certificates. So does the waiting-time logic. A first success is expected after about
$1/(rkN)$ rounds, so with $10^4$--$10^5$ prompts and $k$ of 8--64, LLM RL can find far rarer compositions than our 123
strict theorems allow, if a comparable share of prompts needs them. The numbers do not transfer. Two practical points
follow, with that caveat.
\begin{itemize}
\item \textbf{Evaluation.} Absence at modest $k$ under plain sampling is not absence. Our knockout bases fail at $k=256$,
yet a million plain attempts or a guided evaluation at $k=256$ reveals the composition.
\item \textbf{Data filtering.} Removing a composition from pretraining kept this EI from acquiring it for 16 rounds, but
1\% of records restored it. We did not test the leakage levels a real filter leaves. In the same organism, 16
demonstrations taught a no-\texttt{DN} base excluded middle, so fine-tuning on demonstrations would likely defeat such a
filter.
\end{itemize}

\paragraph{An evidence hierarchy for ``RL created it''.} In increasing strength:
\begin{enumerate}
\item success at equal $k$;
\item a compute-matched base budget, reported as a curve over defensible allocations;
\item RL-free controls with matched training;
\item a certificate against the base at that budget;
\item an intervention on what the base saw.
\end{enumerate}
For our full-exposure seeds the claim shrinks at each step (Table~\ref{tab:budgets}). With the intervention, exposure was
followed by acquisition within 16 rounds, and its absence by none, for one composition and one algorithm. Steps 2 and 3
cost about as much as the RL run they audit, so they are affordable wherever that run was; they are rarely reported.
